\par\Needspace{6\baselineskip}\medskip
\begin{center}\hypertarget{native-166-article120}{}\pdfbookmark[2]{120.}{article-120}120.\end{center}
Quoniam itaque problematis natura reductionem ulteriorem non permittit,
quam ad duas aequationes, duas incognitas mixtim implicantes, rei summa
primo quidem in idonea harum incognitarum \emph{electione} aequationumque
\SourceLeaf{167}\emph{adornatione} versabitur, ut tum $X$ et $Y$ quam simplicissime ab $x,y$ pendeant, tum ex
harum valoribus inventis elementa ipsa quam commodissime demanent: dein
vero circumspiciendum erit, quo pacto incognitarum valores aequationibus satisfacientes per operationes non nimis operosas eruere liceat. Quod si coecis
quasi tentaminibus tantum efficiendum esset, ingens sane ac vix tolerandus
labor requireretur, qualem fere nihilominus saepius susceperunt astronomi,
qui cometarum orbitas per methodum quam indirectam vocant determinaverunt:
magnopere utique in tali negotio labor sublevatur eo, quod in tentaminibus
primis calculi crassiores sufficiunt, donec ad valores approximatos incognitarum
perventum fuerit. Quamprimum vero determinatio approximata iam habetur,
rem tutis semper expeditisque methodis ad finem perducere licebit, quas antequam ulterius progrediamur hic explicavisse iuvabit.

Aequationibus $X=0$, $Y=0$, si pro $x,y$ valores veri ipsi accipiuntur,
ex asse sponte satisfiet: contra si pro $x,y$ valores a veris diversi substituuntur,
$X$ et $Y$ inde valores a $0$ diversos nanciscentur. Quo propius vero illi
ad veros accedunt, eo minores quoque valores ipsarum $X,Y$ emergere debebunt,
quotiesque illorum differentiae a veris perexiguae sunt, supponere licebit,
variationes in valoribus ipsarum $X,Y$ proxime proportionales esse variationi
ipsius $x$, si $y$, vel variationi ipsius $y$, si $x$ non mutetur. Quodsi itaque
valores veri ipsarum $x,y$ resp. designantur per $\xi,\eta$, valores ipsarum $X,Y$ suppositioni
$x=\xi+\lambda$, $y=\eta+\mu$ respondentes per formam $X=\alpha\lambda+\beta\mu$, $Y=\gamma\lambda+\delta\mu$ exhibebuntur, ita ut coëfficientes $\alpha,\beta,\gamma,\delta$ pro constantibus haberi
queant, dum $\lambda$ et $\mu$ perexiguae manent. Hinc concluditur, si pro tribus systematibus
valorum ipsarum $x,y$, a veris parum diversorum, valores respondentes
ipsarum $X,Y$ determinati sint, valores veros ipsarum $x,y$ inde derivari
posse, quatenus quidem suppositionem istam admittere licet. Statuamus
\[
\begin{array}{rlllll}
\text{pro}\quad & x=a, & y=b & \text{fieri}\quad & X=A, & Y=B\\
& x=a', & y=b' && X=A', & Y=B'\\
& x=a'', & y=b'' && X=A'', & Y=B''
\end{array}
\]
habebimusque
\[
\begin{aligned}
A&=\alpha(a-\xi)+\beta(b-\eta),\\
B&=\gamma(a-\xi)+\delta(b-\eta)\\
A'&=\alpha(a'-\xi)+\beta(b'-\eta),\\
B'&=\gamma(a'-\xi)+\delta(b'-\eta)\\
A''&=\alpha(a''-\xi)+\beta(b''-\eta),\\
B''&=\gamma(a''-\xi)+\delta(b''-\eta).
\end{aligned}
\]
\SourceLeaf{168}Hinc fit, eliminatis $\alpha,\beta,\gamma,\delta$
\[
\begin{aligned}
\xi&=\frac{a(A'B''-A''B')+a'(A''B-AB'')+a''(AB'-A'B)}{A'B''-A''B'+A''B-AB''+AB'-A'B}\\[6pt]
\eta&=\frac{b(A'B''-A''B')+b'(A''B-AB'')+b''(AB'-A'B)}{A'B''-A''B'+A''B-AB''+AB'-A'B}
\end{aligned}
\]
sive in forma ad calculum commodiori
\[
\begin{aligned}
\xi&=a+\frac{(a'-a)(A''B-AB'')+(a''-a)(AB'-A'B)}{A'B''-A''B'+A''B-AB''+AB'-A'B}\\[6pt]
\eta&=b+\frac{(b'-b)(A''B-AB'')+(b''-b)(AB'-A'B)}{A'B''-A''B'+A''B-AB''+AB'-A'B}\mathpunct{\cdot}
\end{aligned}
\]
Manifesto quoque in his formulis quantitates $a,b,A,B$ cum $a',b',A',B'$, vel
cum his $a'',b'',A'',B''$ permutare licet.

Ceterum denominator communis omnium harum expressionum, quem etiam
sub formam $(A'-A)(B''-B)-(A''-A)(B'-B)$ ponere licet, fit
\[
=(\alpha\delta-\beta\gamma)\{(a'-a)(b''-b)-(a''-a)(b'-b)\}:
\]
unde patet, $a,a',a'',b,b',b''$ ita accipi debere, ut non fiat
$\dfrac{a''-a}{b''-b}=\dfrac{a'-a}{b'-b}\mathpunct{\raisebox{.45ex}{,}}$
alioquin enim haec methodus haud applicabilis esset, sed pro $\xi$ et $\eta$ valores
fractos suggereret, quorum numeratores et denominatores simul evanescerent.
Simul hinc manifestum est, si forte fiat $\alpha\delta-\beta\gamma=0$, eundem defectum methodi
usum omnino destruere, quomodocunque $a,a',a'',b,b',b''$ accipiantur.
In tali casu pro valoribus ipsius $X$ formam talem supponere oporteret $\alpha\lambda+\beta\mu+\epsilon\lambda\lambda+\zeta\lambda\mu+\theta\mu\mu$, similemque pro valoribus ipsius $Y$, quo facto analysis
methodos praecedenti analogas suppeditaret, e valoribus ipsarum $X,Y$ pro
quatuor systematibus valorum ipsarum $x,y$ computatis harum valores veros
eruendi. Hoc vero modo calculus permolestus evaderet, praetereaque ostendi
potest, in tali casu orbitae determinationem praecisionem necessariam per
ipsius rei naturam non admittere: quod incommodum quum aliter evitari nequeat,
nisi novis observationibus magis idoneis adscitis, huic argumento hic
non immoramur.
